\chapter{Mathematical Foundations}
\label{app:math_foundations}

This appendix provides expanded mathematical background for the classical models
discussed in Section~2.2 and the stochastic framework underlying them. It is
intended as a self-contained reference that deepens the treatment of Brownian
motion, It\^{o}'s Lemma, the Poisson process, and the structural limitations of
GARCH and Merton's jump-diffusion model, without repeating the main exposition.
Throughout this appendix, $t$ denotes calendar time; when referencing
Eq.~(\ref{eq:blackscholes}) from the main text, $t$ plays the role of $\tau$.

\section{Brownian Motion}
\label{app:sec:bm}

Brownian motion is the continuous-time stochastic process that underlies both
the white-noise innovations of GARCH and the diffusion component of the GBM in
Eq.~(\ref{eq:blackscholes}). It was first formally constructed by Wiener (1923)
and is therefore also called the \textit{Wiener process}; the name
\textit{Brownian motion} traces back to the botanist Robert Brown, who in 1827
observed the erratic movement of pollen particles suspended in water, a
phenomenon later explained by Einstein (1905) as the effect of random molecular
collisions.\\

\textbf{Definition.} A stochastic process $\{W_t\}_{t \geq 0}$ defined on a
probability space $(\Omega, \mathcal{F}, \mathbb{P})$ is called a
\textit{standard Brownian motion} if it satisfies the following four properties:
\begin{enumerate}
    \item \textit{Initialisation}: $W_0 = 0$ almost surely.
    \item \textit{Independent increments}: for any $0 \leq t_0 < t_1 < \cdots <
    t_n$, the increments $W_{t_1} - W_{t_0},\, W_{t_2} - W_{t_1},\, \ldots,\,
    W_{t_n} - W_{t_{n-1}}$ are mutually independent random variables.
    \item \textit{Stationary Gaussian increments}: for any $0 \leq s < t$,
    \begin{equation}
        \label{eq:bm_increment}
        W_t - W_s \;\sim\; \mathcal{N}(0,\, t - s).
    \end{equation}
    \item \textit{Continuous paths}: the map $t \mapsto W_t$ is almost surely
    continuous.
\end{enumerate}

These four axioms are minimal and mutually consistent: any two of them, together
with mild regularity, imply the others \cite{oksendal2003stochastic}.\\

\textbf{Intuition: the limit of a random walk.} Brownian motion can be understood
as the continuous-time limit of a symmetric random walk. Partition $[0, T]$ into
$N$ steps of length $\Delta t = T/N$, and at each step $k$ add $+\sqrt{\Delta t}$
or $-\sqrt{\Delta t}$ with equal probability. The position after $k$ steps is
$Z_k = \sum_{i=1}^k \xi_i \sqrt{\Delta t}$ where $\xi_i \in \{-1, +1\}$ with
equal probability. The variance of $Z_k$ is $k \cdot \Delta t$, which remains
finite as $N \to \infty$ provided $\Delta t = T/N$. By the Central Limit
Theorem, $Z_{\lfloor tN \rfloor}$ converges in distribution to $W_t \sim
\mathcal{N}(0, t)$. The $\sqrt{\Delta t}$ scaling is not optional: scaling by
$\Delta t$ would collapse all increments to zero, while scaling by a constant
would produce infinite variance. This is the only scaling that gives a
well-defined non-degenerate continuous-time limit.

\subsection{Key Statistical Properties}

Setting $s = 0$ in property~3, we immediately obtain:
\[
    \mathbb{E}[W_t] = 0, \qquad \operatorname{Var}(W_t) = t, \qquad
    W_t \sim \mathcal{N}(0,\, t).
\]
The variance grows linearly in time: uncertainty accumulates at rate $\sqrt{t}$,
not $t$. This sub-linear growth is the signature of diffusion and contrasts with
a deterministic path (variance 0) or with a purely random variable with no time
structure.\\

\textbf{Covariance and correlation.} For $0 \leq s \leq t$, write
$W_t = W_s + (W_t - W_s)$. Since $W_t - W_s$ is independent of $W_s$ (property~2)
and has mean zero:
\begin{align*}
    \operatorname{Cov}(W_s, W_t)
    &= \operatorname{Cov}\bigl(W_s,\; W_s + (W_t - W_s)\bigr) \\
    &= \operatorname{Var}(W_s) + \underbrace{\operatorname{Cov}(W_s,\,
       W_t - W_s)}_{=\,0}
    \;=\; s \;=\; \min(s, t).
\end{align*}
Consequently $\operatorname{Corr}(W_s, W_t) = \sqrt{s/t}$ for $s \leq t$.
Two observations far apart in time are weakly correlated; the process has no
long-range memory beyond the information contained in its current value, a
property formalised next.\\

\textbf{Markov property.} The future distribution of $W_t$ given all past
information $\mathcal{F}_s = \sigma(W_u : u \leq s)$ depends on the past only
through the current value $W_s$:
\[
    P(W_t \in A \mid \mathcal{F}_s) = P(W_t \in A \mid W_s), \qquad s \leq t.
\]
This follows from independent increments: $W_t - W_s \sim \mathcal{N}(0, t-s)$
is independent of $\mathcal{F}_s$, so knowing the full path up to $s$ adds
nothing beyond knowing $W_s$. Economically, this says that the noise driving
asset prices has no exploitable history: past increments of $W$ carry no
information about future ones.\\

\textbf{Martingale property.} With respect to $\{\mathcal{F}_t\}$, $W_t$ is a
martingale:
\[
    \mathbb{E}[W_t \mid \mathcal{F}_s] = W_s, \qquad s \leq t.
\]
Proof: $\mathbb{E}[W_t \mid \mathcal{F}_s] = \mathbb{E}[W_s + (W_t - W_s)
\mid \mathcal{F}_s] = W_s + \mathbb{E}[W_t - W_s] = W_s$, using independence
and zero mean of the increment. The martingale property encodes the absence of
predictable drift in the noise: knowing the path up to $s$, the best prediction
of $W_t$ is simply $W_s$, with no systematic upward or downward tendency.

\subsection{Quadratic Variation and Non-Differentiability}

A central and counter-intuitive property of BM is that its paths are almost
surely continuous yet nowhere differentiable. This is not a pathological edge
case: it holds with probability one and is intimately connected to the
\textit{quadratic variation} of the process.\\

\textbf{Quadratic variation.} For a partition $\Pi = \{0 = t_0 < t_1 < \cdots
< t_n = t\}$ with mesh $\|\Pi\| = \max_i(t_{i+1} - t_i)$, the quadratic
variation of $W$ over $[0, t]$ is:
\begin{equation}
    \label{eq:quadratic_variation}
    [W]_t \;=\; \lim_{\|\Pi\| \to 0}\; \sum_{i=0}^{n-1}(W_{t_{i+1}} -
    W_{t_i})^2 \;=\; t,
\end{equation}
where the limit holds in $L^2$. To verify this, write $\delta_i = W_{t_{i+1}}
- W_{t_i}$:
\[
    \mathbb{E}\!\left[\sum_i \delta_i^2\right] = \sum_i(t_{i+1} - t_i) = t,
\]
\[
    \operatorname{Var}\!\left[\sum_i \delta_i^2\right]
    = \sum_i \operatorname{Var}(\delta_i^2)
    = \sum_i 2(t_{i+1} - t_i)^2
    \leq 2\|\Pi\|\sum_i(t_{i+1} - t_i)
    = 2\|\Pi\|\, t \;\xrightarrow{\|\Pi\|\to 0}\; 0.
\]
The mean is $t$ and the variance vanishes, so the sum converges in $L^2$ to $t$,
establishing Eq.~(\ref{eq:quadratic_variation}). This result is written
compactly in It\^{o} calculus as the \textit{multiplication rules}:
\begin{equation}
    \label{eq:ito_product_rules}
    (dW_t)^2 = dt, \qquad dW_t\cdot dt = 0, \qquad (dt)^2 = 0.
\end{equation}

For comparison, any deterministic differentiable function $f$ has zero quadratic
variation: its increments are $O(\Delta t)$ so their squares sum to
$O(\|\Pi\|) \to 0$. The non-zero quadratic variation of BM means that it
oscillates so violently on small time scales that the sum of squared increments
is macroscopically significant. Two direct consequences follow:
\begin{itemize}
    \item \textit{Infinite total variation}: $\sum_i |\delta_i| \to \infty$ as
    $\|\Pi\| \to 0$ almost surely. BM paths cannot be approximated by smooth
    curves in the sense of bounded variation.
    \item \textit{Nowhere differentiability}: if $W$ were differentiable at
    some point $t$, then $|\delta_i| = O(\Delta t)$, giving $\sum_i \delta_i^2
    = O(\|\Pi\|) \to 0$, contradicting $[W]_t = t > 0$. Hence no derivative
    exists almost surely at any point.
\end{itemize}

Financially, nowhere differentiability means that instantaneous
``speed'' $dW_t/dt$ does not exist: price changes driven by BM are inherently
irregular at every time scale and carry no predictable directional momentum from
one instant to the next.

\section{It\^{o}'s Lemma and the Black-Scholes Model}
\label{app:sec:ito}

Ordinary calculus provides the chain rule: if $f$ is smooth and $x(t)$ is
differentiable, then $df = f'(x)\,dx$. Because BM paths are not differentiable,
this rule fails when $x$ is driven by a Brownian motion. The extra term in the
quadratic variation, $(dW_t)^2 = dt \neq 0$, means that the second-order
contribution in the Taylor expansion is of order $dt$ and cannot be discarded.
It\^{o}'s Lemma provides the correct generalisation.

\subsection{Statement and Heuristic Proof}

Let $\{X_t\}$ be an It\^{o} process satisfying:
\[
    dX_t = a_t\, dt + b_t\, dW_t,
\]
where $a_t$ (the drift) and $b_t$ (the diffusion coefficient) are adapted
stochastic processes satisfying standard integrability conditions
\cite{oksendal2003stochastic}. Let $f(t, x)$ be a function that is once
continuously differentiable in $t$ and twice continuously differentiable in $x$.
Then:
\begin{equation}
    \label{eq:ito_lemma}
    df(t, X_t) = \left(\frac{\partial f}{\partial t}
    + a_t\, \frac{\partial f}{\partial x}
    + \frac{1}{2}\, b_t^2\,
    \frac{\partial^2 f}{\partial x^2}\right)\! dt
    + b_t\, \frac{\partial f}{\partial x}\, dW_t.
\end{equation}

The term $\frac{1}{2} b_t^2 \frac{\partial^2 f}{\partial x^2}\, dt$ is the
\textbf{It\^{o} correction}. It is absent in ordinary calculus and arises
because the stochastic increment has a non-negligible second-order contribution.
In vector form, for a function $f(t, \mathbf{X}_t)$ of a
multi-dimensional It\^{o} process $d\mathbf{X}_t = \mathbf{a}_t\,dt +
\mathbf{B}_t\,d\mathbf{W}_t$, the correction becomes
$\frac{1}{2}\operatorname{tr}\!\left(\mathbf{B}_t\mathbf{B}_t^{\top}
\nabla^2 f\right)dt$, but the one-dimensional case suffices for the
applications in Chapter~2.\\

\textit{Heuristic derivation.} Expand $f(t + dt, X_{t+dt})$ by Taylor's theorem
to second order in $dX_t$ and first order in $dt$:
\[
    df = \frac{\partial f}{\partial t}\, dt
       + \frac{\partial f}{\partial x}\, dX_t
       + \frac{1}{2}\frac{\partial^2 f}{\partial x^2}(dX_t)^2
       + \cdots
\]
Substitute $dX_t = a_t\, dt + b_t\, dW_t$ and expand $(dX_t)^2$:
\[
    (dX_t)^2 = a_t^2(dt)^2 + 2a_t b_t\, dt\, dW_t + b_t^2(dW_t)^2.
\]
Applying the rules in Eq.~(\ref{eq:ito_product_rules}):
\[
    (dX_t)^2 = b_t^2\, dt.
\]
The cross term $dt\,dW_t$ and the higher-order term $(dt)^2$ vanish. Collecting
the surviving contributions:
\[
    df = \frac{\partial f}{\partial t}\, dt
       + \frac{\partial f}{\partial x}(a_t\, dt + b_t\, dW_t)
       + \frac{1}{2}\frac{\partial^2 f}{\partial x^2}\, b_t^2\, dt,
\]
which rearranges to Eq.~(\ref{eq:ito_lemma}). The rigorous proof replaces
this Taylor expansion argument with an $L^2$ limit over sequences of
partitions \cite{oksendal2003stochastic}; the intuition is identical.

\subsection{Proof of the Black-Scholes Price Equation}

Chapter~2 states that applying It\^{o}'s Lemma to $\log S_t$ under GBM
(Eq.~\ref{eq:blackscholes}) yields the log-normal price formula. We derive
this in full here. The GBM is the It\^{o} process with drift coefficient
$a_t = \mu S_t$ and diffusion coefficient $b_t = \nu S_t$:
\[
    dS_t = \mu\, S_t\, dt + \nu\, S_t\, dW_t.
\]
Apply It\^{o}'s Lemma (Eq.~\ref{eq:ito_lemma}) to $f(S_t) = \log S_t$:
\[
    \frac{\partial f}{\partial t} = 0, \qquad
    \frac{\partial f}{\partial S} = \frac{1}{S_t}, \qquad
    \frac{\partial^2 f}{\partial S^2} = -\frac{1}{S_t^2}.
\]
Substituting $a_t = \mu S_t$ and $b_t = \nu S_t$:
\begin{align*}
    d(\log S_t)
    &= \left(0 + \mu S_t \cdot \frac{1}{S_t}
       + \frac{1}{2}(\nu S_t)^2 \cdot \left(-\frac{1}{S_t^2}\right)\right)\! dt
       + \nu S_t \cdot \frac{1}{S_t}\, dW_t \\
    &= \left(\mu - \frac{1}{2}\nu^2\right)\! dt + \nu\, dW_t.
\end{align*}
This is a Brownian motion with constant drift $\mu - \frac{1}{2}\nu^2$ and
constant diffusion coefficient $\nu$. Integrating over $[0, t]$:
\[
    \log S_t - \log S_0
    = \left(\mu - \frac{1}{2}\nu^2\right)t + \nu W_t,
\]
and exponentiating:
\begin{equation}
    \label{eq:gbm_solution_app}
    S_t = S_0 \exp\!\left[\left(\mu - \frac{1}{2}\nu^2\right)t
          + \nu W_t\right].
\end{equation}
Since $W_t \sim \mathcal{N}(0, t)$, the log-return over $[0, t]$ satisfies:
\begin{equation}
    \label{eq:lognormal_returns_app}
    \log\!\left(\frac{S_t}{S_0}\right)
    \;\sim\; \mathcal{N}\!\left(\left(\mu - \frac{1}{2}\nu^2\right)t,\;
    \nu^2 t\right).
\end{equation}
Log-returns are normally distributed with mean proportional to $t$ and variance
proportional to $t$, confirming the statement in Section~2.2 and the foundation
of the Black--Scholes option pricing model \cite{black1973pricing}.\\

\textbf{Interpreting the It\^{o} correction.} The term $-\frac{1}{2}\nu^2$ is
not an approximation error: it is the exact price of the convexity of the
exponential function. Because $x \mapsto e^x$ is convex, Jensen's inequality
gives $\mathbb{E}[e^X] \geq e^{\mathbb{E}[X]}$ for any random variable $X$. For
$X = (\mu - \frac{1}{2}\nu^2)t + \nu W_t$, we have $\mathbb{E}[X] = (\mu -
\frac{1}{2}\nu^2)t$ and $\mathbb{E}[e^X] = e^{\mu t}$; indeed:
\[
    \mathbb{E}[S_t] = S_0\, e^{(\mu - \frac{1}{2}\nu^2)t}\,
    \mathbb{E}\!\left[e^{\nu W_t}\right]
    = S_0\, e^{(\mu - \frac{1}{2}\nu^2)t}\cdot e^{\frac{1}{2}\nu^2 t}
    = S_0\, e^{\mu t},
\]
where $\mathbb{E}[e^{\nu W_t}] = e^{\frac{1}{2}\nu^2 t}$ is the moment
generating function of the Gaussian. The correction ensures that the arithmetic
expected return of prices is $\mu$ (as specified in the SDE), while the
geometric mean return is the lower quantity $\mu - \frac{1}{2}\nu^2$. This
gap, $\frac{1}{2}\nu^2$, is called \textit{variance drag}: higher volatility
reduces the realised compound return even holding the instantaneous drift fixed.
Without the It\^{o} correction, applying the classical chain rule would give
$d(\log S_t) = \mu\, dt + \nu\, dW_t$, implying $\log(S_t/S_0) \sim
\mathcal{N}(\mu t, \nu^2 t)$ and therefore $\mathbb{E}[S_t] = S_0
e^{\mu t + \frac{1}{2}\nu^2 t}$, which contradicts $\mathbb{E}[S_t] = S_0
e^{\mu t}$ derived directly from the SDE.

\section{Poisson Process and Merton's Jump-Diffusion Model}
\label{app:sec:poisson_merton}

The Poisson process is the canonical model for events that arrive randomly,
infrequently, and with no predictable timing — in the financial context, sudden
discontinuous jumps in asset prices caused by earnings surprises, credit events,
regulatory changes, or macroeconomic shocks. It complements Brownian motion by
introducing discontinuous sample paths, allowing the model to generate
large price moves instantaneously rather than through accumulated small
increments.

\subsection{The Poisson Process}

\textbf{Definition.} A stochastic process $\{P_t\}_{t \geq 0}$ taking values in
$\mathbb{N}_0 = \{0, 1, 2, \ldots\}$ is a \textit{Poisson process with intensity}
$\lambda > 0$ if:
\begin{enumerate}
    \item $P_0 = 0$ almost surely.
    \item Independent increments: for any $0 \leq t_0 < t_1 < \cdots < t_n$,
    the increments $P_{t_1} - P_{t_0},\, \ldots,\, P_{t_n} - P_{t_{n-1}}$ are
    mutually independent.
    \item Stationary Poisson increments: for any $0 \leq s < t$,
    \begin{equation}
        \label{eq:poisson_increment}
        P_t - P_s \;\sim\; \operatorname{Poisson}(\lambda(t - s)),
    \end{equation}
    i.e., $P(P_t - P_s = k) = \frac{e^{-\lambda(t-s)}[\lambda(t-s)]^k}{k!}$
    for $k \in \mathbb{N}_0$.
    \item At most one event per instant: $P(P_{t+h} - P_t \geq 2) = o(h)$
    as $h \to 0^+$.
\end{enumerate}

The process $P_t$ counts the number of jumps that have occurred by time $t$.
From Eq.~(\ref{eq:poisson_increment}):
\[
    \mathbb{E}[P_t] = \lambda t, \qquad \operatorname{Var}(P_t) = \lambda t.
\]
The equality of mean and variance is the defining property of the Poisson
distribution: the intensity $\lambda$ simultaneously controls the average
jump rate and the variability of the jump count. A small $\lambda$ means
rare, widely-spaced jumps; a large $\lambda$ means frequent ones.\\

\textbf{Inter-arrival times.} The waiting times between consecutive jumps
$\tau_1, \tau_2, \ldots$ are independent and exponentially distributed with
mean $1/\lambda$:
\[
    \tau_k \;\sim\; \operatorname{Exp}(\lambda), \qquad
    P(\tau_k > s) = e^{-\lambda s}.
\]
The memoryless property of the exponential distribution mirrors the independent
increments of $P_t$: the time until the next jump does not depend on how long
one has already waited. This is a deliberate modelling choice; in reality, jump
arrivals in markets are not memoryless (see Section~\ref{app:sec:limitations}).\\

\textbf{Martingale structure.} The \textit{compensated Poisson process}
$\widetilde{P}_t = P_t - \lambda t$ is a martingale with zero mean and
$\operatorname{Var}(\widetilde{P}_t) = \lambda t$. The deterministic term
$\lambda t$ is the \textit{compensator}: it removes the upward drift of $P_t$,
playing the same role for the Poisson process that zero mean plays for $W_t$.

\subsection{Merton's Jump-Diffusion Model}

Merton \cite{merton1976option} extended the GBM (Eq.~\ref{eq:blackscholes}) by
superimposing a compound Poisson jump component onto the continuous diffusion:
\begin{equation}
    \label{eq:merton_sde_app}
    dS_t = \mu\, S_t\, dt + \nu\, S_t\, dW_t + S_{t^-}(J - 1)\, dP_t.
\end{equation}
The three terms have distinct roles:
\begin{itemize}
    \item $\mu\, S_t\, dt$: the continuous expected-return drift, as in GBM.
    \item $\nu\, S_t\, dW_t$: the diffusion noise, modelling continuous
    day-to-day fluctuations with constant volatility $\nu > 0$.
    \item $S_{t^-}(J - 1)\, dP_t$: the jump term. Here $S_{t^-}$ denotes
    the price immediately \textit{before} a possible jump at time $t$,
    $P_t$ is a Poisson process with intensity $\lambda$ independent of $W_t$,
    and $J > 0$ is a random jump multiplier, i.i.d.\ across jump events and
    independent of both $W_t$ and $P_t$. When $dP_t = 1$, the price jumps
    from $S_{t^-}$ to $S_{t^-} \cdot J$; when $dP_t = 0$, the jump term
    vanishes.
\end{itemize}

\textbf{Distribution of the jump multiplier.} As discussed in Section~2.2, the
log-return distribution under Merton's model is a Poisson-weighted mixture of
Gaussians. This result requires $\log J$ to be normally distributed; consistent
with Merton's (1976) original formulation, $J$ is therefore \textit{log-normally}
distributed:
\begin{equation}
    \label{eq:jump_lognormal_app}
    \log J \;\sim\; \mathcal{N}(\mu_j,\, \sigma_J^2).
\end{equation}
The parameter $\mu_j$ is the mean log-jump size: $\mu_j > 0$ corresponds to
jumps that are on average upward, $\mu_j < 0$ to downward jumps, and $\mu_j = 0$
to symmetric jumps. The parameter $\sigma_J^2$ controls the dispersion of jump
sizes; larger $\sigma_J$ produces a wider range of jump magnitudes and heavier
tails. The log-normal specification also guarantees $J > 0$, ensuring that
prices remain positive after any jump.

\subsection{Log-Return Distribution under Merton's Model}

Applying the jump-diffusion analogue of It\^{o}'s Lemma to $f(S_t) = \log S_t$
under Eq.~(\ref{eq:merton_sde_app}), the continuous part is handled exactly as
in Appendix~\ref{app:sec:ito}, while each jump at time $t$ contributes
$\log(J)$ to the log-price. The resulting log-price dynamics are:
\[
    d(\log S_t) = \left(\mu - \frac{1}{2}\nu^2\right)\! dt
                + \nu\, dW_t + \log J\, dP_t.
\]
Integrating over $[0, t]$, conditional on exactly $n$ jumps occurring in
$[0, t]$ (i.e., $P_t = n$), the log-return is a sum of the continuous
Gaussian component and $n$ independent log-normal jump contributions:
\[
    \log\!\left(\frac{S_t}{S_0}\right) \;\Bigg|\; P_t = n
    \;\sim\; \mathcal{N}\!\left(\left(\mu - \tfrac{1}{2}\nu^2\right)t
    + n\mu_j,\;\; \nu^2 t + n\sigma_J^2\right).
\]
The unconditional log-return density is obtained by summing over all
possible jump counts, weighted by their Poisson probabilities
$p_n = e^{-\lambda t}(\lambda t)^n / n!$:
\begin{equation}
    \label{eq:merton_return_dist}
    p\!\left(\log\tfrac{S_t}{S_0} = x\right)
    = \sum_{n=0}^{\infty}
    \frac{e^{-\lambda t}(\lambda t)^n}{n!}\;
    \phi\!\left(x;\;
    \left(\mu - \tfrac{1}{2}\nu^2\right)t + n\mu_j,\;
    \nu^2 t + n\sigma_J^2\right),
\end{equation}
where $\phi(\cdot;\, m, v)$ denotes the Gaussian density with mean $m$ and
variance $v$. This is the Poisson-weighted mixture of Gaussians stated in
Section~2.2.\\

\textbf{Intuition.} Each term in the sum corresponds to a scenario with exactly
$n$ jumps. The $n = 0$ term (probability $e^{-\lambda t}$, close to one when
$\lambda$ is small) is the pure GBM return, which is Gaussian. Each additional
jump shifts the component mean by $\mu_j$ and inflates the component variance
by $\sigma_J^2$, spreading the overall distribution. Three properties emerge:
\begin{itemize}
    \item \textit{Leptokurtosis}: the mixture has heavier tails than a single
    Gaussian because the variance-inflated components for $n \geq 1$ place more
    probability mass on large deviations.
    \item \textit{Skewness}: when $\mu_j \neq 0$, the jump mean shifts all
    components with $n \geq 1$, creating asymmetry in the unconditional
    distribution.
    \item \textit{Multi-modality} (for large $|\mu_j|$): if the jump mean is
    large relative to the component widths, distinct peaks can emerge
    corresponding to the zero-jump and one-jump components.
\end{itemize}
These properties represent a genuine improvement over pure GBM in capturing
the stylized facts of fat tails and occasional large price moves that are
relevant to financial time series (Section~2.1).

\section{Further Limitations of GARCH and Merton's Model}
\label{app:sec:limitations}

Section~2.2 identifies the core limitation shared by both models: parametric
distributional assumptions restrict their expressivity, and both require
re-estimation whenever the data-generating regime changes. This section
elaborates on the structural limitations of each model individually, expanding
on why they remain insufficient baselines for capturing all the stylized facts
described in Section~2.1.

\subsection{Further Limitations of GARCH}

\textbf{Symmetric shock response and the leverage effect.} The GARCH(1,1)
variance equation $\sigma_t^2 = \omega + \alpha\epsilon_{t-1}^2 +
\beta\sigma_{t-1}^2$ depends on $\epsilon_{t-1}^2$, which is symmetric in the
sign of the innovation: a negative shock $\epsilon_{t-1} = -c$ and a positive
shock $\epsilon_{t-1} = +c$ of the same magnitude produce identical increases in
$\sigma_t^2$. This contradicts the leverage effect (Section~2.1), which
documents that negative returns are followed by disproportionately higher
future volatility. By construction, GARCH(1,1) is unable to produce this
asymmetry. Extensions such as GJR-GARCH and EGARCH introduce asymmetric terms,
but at the cost of additional parameters and reduced parsimony.\\

\textbf{Tail inadequacy.} Even when innovations are drawn from a Student-$t$ or
Generalised Error Distribution rather than a Gaussian, the tail behaviour of the
GARCH conditional distribution remains insufficiently heavy for many financial
series. Section~2.1 documents power-law tails with exponents $\alpha \in [3,5]$,
implying that returns have finite variance but potentially infinite or unstable
fourth moments. Standard GARCH with Gaussian innovations has finite unconditional
moments of all orders, and while heavier-tailed innovations extend the tail
decay, precisely matching an empirical power-law over the full range of quantiles
requires distributional choices that are not robust across assets or time
periods.\\

\textbf{Integrated GARCH and non-stationarity.} The stationarity condition
$\alpha + \beta < 1$ ensures that the unconditional variance
$\bar{\sigma}^2 = \omega / (1 - \alpha - \beta)$ is finite. In practice,
maximum likelihood estimation on financial returns frequently yields
$\hat\alpha + \hat\beta \approx 1$, placing the model on the boundary of
non-stationarity (IGARCH), or even outside it on short samples. In the IGARCH
case, the unconditional variance is infinite, making the model formally
inconsistent with any stationary long-run distribution, yet short-horizon
conditional variance forecasts remain well-defined. This boundary behaviour
is not cleanly handled within the GARCH(1,1) framework and complicates both
inference and the interpretation of estimated parameters.\\

\textbf{Short-memory volatility autocorrelation.} In GARCH(1,1), the
autocorrelation of squared returns decays exponentially at rate
$(\alpha + \beta)^k$ for lag $k$. Empirically, as noted in Section~2.1,
volatility autocorrelation decays hyperbolically,
$\operatorname{Corr}(|r_t|, |r_{t+k}|) \sim k^{-\beta}$ with
$\beta \in [0.1, 0.5]$ \cite{cont2001empirical}, characteristic of a
long-memory process. This mismatch is not a minor calibration issue: exponential
decay systematically underestimates volatility persistence at long horizons, even
when $\alpha + \beta$ is close to one.\\

\textbf{Absence of jumps.} GARCH generates smooth conditional variance paths in
discrete time. It cannot produce discontinuous price gaps or isolated large moves
of the kind associated with earnings surprises, default events, or flash crashes.
High-volatility episodes can cause large returns, but the returns must accumulate
step by step; an instantaneous price discontinuity is structurally excluded.\\

\textbf{Multivariate extension.} Extending GARCH to $K$ assets requires modelling
the full $K \times K$ conditional covariance matrix. Models such as DCC-GARCH
parametrize the correlation dynamics, but the number of free parameters still
grows as $O(K^2)$, making estimation numerically fragile for moderate $K$ and
effectively infeasible for the high-dimensional cross-sections that are typical
of portfolio applications. This is a fundamental obstacle for joint modelling of
multi-asset systems, unlike the deep generative models discussed in Sections~2.3
and~2.4.\\

\textbf{Re-estimation requirement.} GARCH parameters are estimated on a
historical window and are fixed at prediction time. If the market transitions
between volatility regimes — for instance from a low-volatility expansion to a
financial crisis — the estimated parameters become stale. The model must then be
re-fitted from scratch on the new regime's data, a process that requires
identifying the regime transition, accumulating sufficient data, and repeating
estimation. Real-time adaptation is not natively supported.

\subsection{Further Limitations of Merton's Model}

\textbf{Constant jump intensity.} The Poisson intensity $\lambda$ in
Eq.~(\ref{eq:merton_sde_app}) is a fixed constant. In reality, jump activity is
time-varying and clustered: large price moves tend to arrive in bursts during
periods of market stress and are rare during calm periods. A constant $\lambda$
cannot reproduce this self-reinforcing jump clustering; capturing it would
require an intensity that depends on the history of past jumps or on latent
market state.\\

\textbf{Tail inadequacy despite fat tails.} Although the Poisson mixture
(Eq.~\ref{eq:merton_return_dist}) generates heavier tails than a single
Gaussian, all moments of the log-return distribution remain finite.
The $n$-th mixture component has variance $\nu^2 t + n\sigma_J^2$, which
grows linearly in $n$, but the Poisson weights $e^{-\lambda t}(\lambda t)^n/n!$
decay super-exponentially in $n$. Consequently, the unconditional tail of
Eq.~(\ref{eq:merton_return_dist}) decays faster than any power law, which is
inconsistent with the empirically observed tail exponents $\alpha \in [3,5]$
\cite{cont2001empirical}. The log-Gaussian jump size assumption therefore
produces excess kurtosis and wide tails in practice, but cannot fully reproduce
the power-law tail decay of real return distributions. Replacing the Gaussian
jump distribution with a heavier-tailed alternative (e.g., a stable or
double-exponential distribution) relaxes this constraint but either introduces
infinite variance or makes the likelihood function intractable.\\

\textbf{No volatility clustering.} The diffusion component $\nu\, S_t\, dW_t$
has constant volatility $\nu$. Jumps introduce excess kurtosis and fat tails,
but they do not produce persistent time-varying conditional variance: between
jump events the model reverts to a constant-volatility GBM, and the occurrence
of a jump does not alter the volatility of subsequent diffusion increments.
Volatility clustering — extended periods of elevated or depressed volatility —
is therefore structurally absent in Merton's model. Combining the jump component
with a stochastic volatility diffusion (as in the Bates extension) partially
addresses this, but requires additional parameters and substantially complicates
estimation.\\

\textbf{No leverage effect.} Both the diffusion noise $\nu\, dW_t$ and the jump
multiplier $J$ are specified symmetrically with respect to the sign of the
return: neither component introduces the asymmetric response whereby negative
returns are followed by higher future volatility more strongly than positive
returns of the same magnitude. Merton's model therefore cannot reproduce the
negative lead-lag correlation $L(k) < 0$ defined in Section~2.1.\\

\textbf{Independence of jumps and diffusion.} The model assumes that the Poisson
jump process and the Brownian diffusion are independent: jumps are no more likely
during periods of elevated diffusion volatility than during calm ones. This is at
odds with empirical evidence that extreme returns cluster around periods of
generally elevated volatility, suggesting a positive dependence between the two
components that the model cannot capture.\\

\textbf{Calibration difficulty.} The likelihood of Merton's model is the infinite
sum in Eq.~(\ref{eq:merton_return_dist}), truncated in practice at a finite
number of terms. This likelihood surface has multiple local optima, and the jump
parameters $(\lambda, \mu_j, \sigma_J^2)$ are poorly identified unless the sample
is long and contains clearly distinguishable jump events. With limited data, a
small number of outliers can dominate the estimation of $\lambda$ and $\sigma_J^2$,
leading to unstable parameter estimates. As with GARCH, all parameters are fixed
after calibration and must be re-estimated if the market regime changes.\\

\textbf{Fixed parametric jump distribution.} More broadly, any fixed
specification of the jump size distribution imposes a parametric constraint on
the shape of the return tails. If the true jump distribution is complex,
multi-modal, or time-varying, a fixed log-Gaussian form will be misspecified in
a way that cannot be corrected by changing $\mu_j$ or $\sigma_J^2$ alone. This
is the fundamental tension between the analytical tractability of Merton's model
and the empirical complexity of real markets, and it motivates the use of
non-parametric, data-driven generative models such as the one developed in this
thesis.
